\documentclass[main.tex]{subfiles}
\begin{document}
\chapter{Foliations in Positive Characteristic and $p$-Curvature} \label{chapter:foliations}


In this chapter we review the definition and basic properties of foliations in positive characteristic. Since there is no reference known to the author containing all these statements in the precise form given here (e.g.\ for general foliations as opposed to flat connections), we spell out some of the non-commutative algebra involved in working with $p$-curvatures.
\par 
First we make a few remarks about conventions. Let $p$ be a prime number. For a scheme $X$ of pure characteristic $p$, let $\Frob_X : X \to X$ denote the absolute Frobenius, the morphism given by the identity on the underlying topological space and with $\Frob_X^{\#} : \cO_X \to \cO_X$ given by $f \mapsto f^p$. Let $f : X \to S$ be a morphism of schemes in characteristic $p$, then the diagram,
\begin{center}
\begin{tikzcd}[column sep = large, row sep = large]
X \arrow[r, "\Frob_X"] \arrow[d, "f"] & X \arrow[d, "f"]
\\
S \arrow[r, "\Frob_S"] & S 
\end{tikzcd}
\end{center}  
The fiber product $X^{(p/S)} := X \times_{S, \Frob_S} S$ is called the \textit{Frobenius twist} of $X$ relative to $S$. This diagram implies the existence of a morphism $F_{X/S} : X \to X^{(p/S)}$ via the universality of the pullback square
\begin{center}
\begin{tikzcd}[column sep = large, row sep = large]
X \arrow[r, "F_{X/S}"] \arrow[rd] \arrow[rr, bend left, "\Frob_X"] & X^{(p/S)} \arrow[d] \pullback \arrow[r] & X \arrow[d]
\\
& S \arrow[r, "\Frob_S"] & S
\end{tikzcd}
\end{center}
which we call the \textit{relative Frobenius} of $X$ over $S$. When there is no ambiguity about the base scheme, such as when $X$ lives over a ground field $k$, we will write $F_{X}$ for $F_{X/k}$. 
\par 
On a (possibly singular) $S$-scheme, we define the \textit{relative tangent sheaf} 
\[ \T_{X/S} := \Omega_{X/S}^{\vee} = \shDer_{\cO_S}(\cO_X, \cO_X) \]
as the dual of the sheaf of K\"{a}hler differentials which equals the sheaf of $S$-derivations on $X$. This is a sheaf of $\cO_S$-linear Lie algebras via the Lie bracket of derivations. When $X$ is a variety over a field $k$ we will write $\Omega_{X}$ for $\Omega_{X/k}$ and likewise $\T_{X}$ for $\T_{X/k}$. Note that the tangent sheaf enjoys much worse functorial properties than the cotangent sheaf. If $f : X \to Y$ is a morphism of $S$-schemes with $Y$ smooth over $S$ then dualizing the pullback map
\[ f^* \Omega_{Y/S} \to \Omega_{X/S} \]
we obtain a map called the \textit{differential}
\[ \d{f} : \T_{X/S} \to f^* \T_{Y/S}. \]
Since the map goes the wrong way, we do not obtain a map of sections of the tangent sheaf. Moreover, if $\Omega_{Y/S}$ is not locally free, the differential map is not well-defined since duals do not commute with pullbacks. In general we only have the following diagram of $\cO_X$-modules,
\begin{center}
\begin{tikzcd}
\T_{X/S} \arrow[r, "\d{f}"] & (f^* \Omega_{Y/S})^{\vee}
\\
& f^* \T_{Y/S} \arrow[u] 
\end{tikzcd}
\end{center}
Therefore, care must be taken when making functoriality arguments in the singular setting. In the following, we usually assume the target is smooth for simplicity.

\section{Foliations}

\begin{defn}
Let $X$ be a normal variety. A \textit{foliation} $\F$ is a coherent subsheaf $\F \subset \T_X$ of the tangent sheaf $\T_X$ of $X$ which is
\begin{enumerate}
    \item \textit{involutive}: $[\F, \F] \subset \F$ meaning $\F$ is closed under the Lie bracket
    \item \textit{saturated}: $\T_X / \F$ is torsion free.
\end{enumerate}
\end{defn}

We will almost always only consider foliations on smooth varieties.
In this text, we will only consider \textit{algebraic} foliations where the objects in the above definition lie in the algebraic category and the sheaves are coherent sheaves of $\cO_X$-modules on finite type schemes. 
One can also consider \textit{analytic} or \textit{smooth} foliations taking subsheaves of the tangent bundle in those corresponding categories. Over $\CC$, the involutive condition ensures that $\F$ is locally \textit{integrable}. Integrability means that through any smooth point of $X$ where $\F$ is a sub-bundle, there is a unique analytic submanifold ``integrating'' $\F$ in the sense that its tangent space everywhere agrees with the subspaces $\F \subset \T_X$. However, in a small analytic neighborhood of such a point, the leaves look like fibers of a smooth (submersive) morphism. This motivates the following definition. 

\begin{defn}
The \textit{singular locus} of a foliation $\F \subset \T_X$ denoted as $\Sing(\F)$ is the union of $X^{\sing}$ and the locus where $\F \subset \T_X$ is not a sub-bundle. 
\end{defn}

The saturation condition implies that a foliation on a normal variety is determined by its restriction to any dense open set. Hence algebraic foliations are of a birational nature and often inform the birational study of $X$. This, along with the following functoriality result, implies that foliations can be pulled back along rational maps, confirming their birational nature.

\begin{defn} \label{defn:pullback_foliation}
Let $f : X \to Y$ be a morphism of varieties and $\F \subset \T_{Y}$ a foliation on $Y$. Then we define the \textit{pullback foliation} $f^+ \F \subset \T_X$ as the saturation of the image of $\G \to \T_Y$ arising in the pullback square
\begin{center}
\begin{tikzcd}
\G \pullback \arrow[r] \arrow[d] & \T_{X} \arrow[d, "\d{f}"]
\\
(f^* \F^{\vee})^{\vee} \arrow[r] & (f^* \Omega_{Y})^{\vee}
\end{tikzcd}
\end{center} 
\end{defn}

If $f$ is generically smooth, this operation preserves the corank, $\rank{\T_X} - \rank{\F}$, of the foliation. 
\par 
Often these analytic leaves are very far from being algebraic as they have wild global behavior. For example, they can be dense in the analytic topology and hence only form \textit{immersed} submanifolds (when restricted away from the singular locus of the foliation). It will also be important to study \textit{algebraic leaves} which are subvarieties tangent to the foliations. This is exactly the special case when some analytic leaf is Zariski closed. 

\begin{defn}
Let $\F \subset \T_X$ be a foliation on $X$. We say a subvariety $Z \embed X$ is \textit{tangent} to $\F$  if the map $\T_Z \to \T_X |_Z$ factors through $\F_Z \subset \T_X|_Z$, the saturation of $(\F|_Z)^{\vee \vee} \embed \T_X|_Z$. If additionally $\T_Z \to \T_X|_Z$ is identified with $\F_Z$ we say that $Z$ \textit{is a leaf} of $\F$.
\end{defn}

Note that the property of being tangent depends only on behavior along a dense open set (e.g.\ along the smooth locus of $Z$). 

\section{$p$-Curvature}

In positive characteristic, there is a new invariant of a foliation called the $p$-curvature. To motivate its definition, we take a moment to consider the ordinary curvature of a subsheaf $\F \subset \T_X$. The \textit{curvature} is the obstruction to $\F$ being a sheaf of sub Lie algebras of $\T_X$ i.e.\ it measures the failure of the involutivity condition $[\F, \F] \subset \F$. In analogy with connections, we sometimes call this condition \textit{flatness}. 

\begin{defn}
If $\F$ is not involutive, then there is a $\cO_X$-linear map called the \textit{curvature} or \textit{residual Lie bracket} witnessing the failure of involutivity. It is defined as the map of $\cO_X$-modules,
\[ \ell_{\F_X} : \bigwedge\nolimits^{\! 2} \F \to \T_X / \F \quad \text{ via } \quad \ell_{\F_X} : X \wedge Y \mapsto [X,Y] \mod \F. \]
The map $\ell_{\F_X}$ is well-defined and $\cO_X$-linear because for any $f \in \cO_X$
\[ f [X,Y] = [f X, Y] + Y(f) \cdot X = [X, f Y] - X(f) \cdot Y \]
and therefore
\[ f [X, Y] \equiv [f X, Y] \equiv [X, f Y] \mod \F \]
because if $X, Y \in \F$ then so are $Y(f) \cdot X$ and $X(f) \cdot Y$ since $\F$ is a $\cO_X$-module.
\end{defn}

In characteristic $p > 0$, the tangent bundle $\T_X$ is not just a Lie algebra, it is a $p$-Lie algebra meaning there is a Frobenius semi-linear operation $(-)^{p}$ commuting with the Lie-bracket defined by $\partial \mapsto \partial^p$ where $\partial^p$ is the differential operator given by iteratively applying $\partial$. 

\begin{lemma}
Let $A$ be an algebra over a field $k$ of characteristic $p$ and $\partial : A \to A$ a $k$-derivation. Then the differential operator $\partial^p$ given by
\[ x \mapsto \underbrace{\partial \cdots \partial}_{p \text{ times}} x \]
is a $k$-derivation.
\end{lemma}

\begin{proof}
It is clear that $\partial^p$ is $k$-linear. The Leibniz rule implies that
\[ \partial^n(fg) = \sum_{k = 0}^n { n \choose k} \partial^k(f) \partial^{n-k}(g). \]
Taking $n = p$ and using the divisibility of binomial coefficients gives
\[ \partial^p(fg) = \partial^p(f) g + f \partial^p(g). \]
\end{proof}

The $p$-curvature is the obstruction to $\F \subset \T_X$ being a subsheaf of $p$-Lie algebras given that it is a Lie algebra subsheaf. 


\begin{defn}
Given an involutive subsheaf $\F \subset \T_X$, the $p$-\textit{curvature} map is the map
\[ \psi_{\F,p} : \Frob^* \F \to \T_X / \F \quad  \text{ via } \quad \psi_{\F,p} : \partial \mapsto \partial^p \mod \F. \]
measuring the failure of $\F$ to be $p$-\textit{closed} i.e.\ $\F^p \subset \F$ meaning $\F \subset \T_X$ is a sub $p$-Lie algebra.
\end{defn}

The $p$-curvature turns out to be an $\cO_X$-linear map of $\cO_X$-modules. The proof of this fact is surprisingly tricky. The easiest route is via the following Jacobson relations and Deligne's formula for $p$-powers in noncommutative rings. By a \textit{Lie polynomial} we mean a noncommutative polynomial formed from addition and compositions of commutators.

\begin{lemma}[Jacobson, Deligne] (cf. \cite[Theorem~5.3]{katz:nilpotent_connections})
Let $R$ be a (possibly noncommutative) associative ring of characteristic $p$. Then,
\begin{enumerate}
\item there are universal Lie polynomials $s_i(x,y) \in \FF_p \left< x, y \right>$ such that for all $a,b \in R$
\[ (a + b)^p = a^p + b^p + \sum_{i = 1}^{p-1} s_i(a,b) \]
where $s_i$ is defined by the universal relation in $\FF_p \left< x, y \right>[t]$,
\[ (\ad_{a t + b})^{p-1}(a) = \sum_{i = 1}^{p-1} i \, s_i(a,b) \, t^{i-1} \]
\item for all $g, \theta \in R$ such that $\{ \ad_{\theta}^n(g^m) \}_{n,m \ge 0}$ commute with each other,
\[ (g \theta)^p = g^p \theta^p + g \cdot \ad_{\theta}^{p-1}(g^{p-1}) \theta. \]
\end{enumerate}
\end{lemma}

\begin{rmk}
Since they arise from applying $\ad_{x t + y}^{p-1}$ to $x$, the polynomials $s_i(x,y)$ lie in $\Lie(x,y)_{p-1}$, the $(p-1)^{\text{th}}$-stage of the lower central series. Recall that for a Lie algebra $\g$ the lower central series is the sequence of Lie subalgebras: $\g_{i+1} = [\g_i, \g]$ setting $\g_0 = \g$. 
\end{rmk}


\begin{example}
For example, if $p = 2$ then $(a + b)^2 = a^2 + ab + ba + b^2$ so $s_1(a,b) = ab + ba = [a,b]$ because the characteristic is 2.
\end{example}

\begin{lemma}
If $\F$ is involutive, the $p$-curvature map $\psi_{\F,p}$ is an $\cO_X$-linear map of $\cO_X$-modules.
\end{lemma}

\begin{proof}
To show $\psi_{\F,p}(\partial_1 + \partial_2) = \psi_p(\partial_1) + \psi_p(\partial_2)$ we use the relation
\[ (\partial_1 + \partial_2)^p = \partial_1^p + \partial_2^p + \sum_{i = 1}^{p-1} s_i(\partial_1, \partial_2) \]
where the $s_i$ are Lie polynomials so $s_i(\partial_1, \partial_2) \in \F$ because $\F$ is closed under Lie bracket.
We further need to show that $\psi_{\F,p}(f \partial) = f^p \psi_p(\partial)$ so that it becomes $\cO_X$-linear viewed as a map on $\Frob_X^* \F$. Indeed, since $\ad_{\partial}^n(f^m) = \partial^n(f^m)$ is a scalar, these commute in the ring of differential operators. Hence by Jacobson's lemma,
\[ (f \partial)^p = f^p \partial^p + f \ad_{\partial}^{p-1}(f^{p-1}) \partial \]
and because $\ad_{\partial}^{p-1}(f^{p-1}) = \partial^{p-1}(f^{p-1}) \in \cO_X$, the second term lies in $\F$.
\end{proof}

\subsection{Compatibility}

It will be important to know that $p$-curvature is compatible with pullbacks. One of the main applications of the $p$-curvature map is that it imposes a restriction on the algebraic leaves of a foliation. When $\F$ is not involutive, the residual Lie bracket map $\ell$ gives a restriction on the locus of integrable subvarieties.

\begin{prop}
Let $X$ be a smooth variety and $f : Z \to X$ a morphism \textit{invariant} under a subsheaf $\F \subset \T_X$ in the sense that the tangent map factors as
\[ \d{f} : \T_Z \to f^* \F \to f^* \T_X \]
then the composition of $\T_Z \to f^* \F$ with the pullback of the curvature,
\[ \bigwedge\nolimits^{\! 2} \T_Z \to \bigwedge\nolimits^{\! 2} f^* \F \xrightarrow{f^* \ell_{\F}} f^* (\T_X / \F)  \]
is identically zero.  
\end{prop}

\begin{proof}
This follows immediately from the following Proposition~\ref{prop:curvature_pullback_compatibility} expressing the compatibility of curvature with pullback applied to the map $f : (Z, \T_Z) \to (\T_X, \F)$ and using that the subsheaf $\T_Z \subset \T_Z$ trivially has zero curvature.
\end{proof}

Of course, if $Z$ is a curve, the curvature of $\F$ imposes no restriction since $\bigwedge\nolimits^{\! 2} \T_Z = 0$. However, this is a nontrivial obstruction to finding invariant subvarieties (or immersed submanifolds in the analytic setting) of dimension $> 1$. 

\begin{prop} \label{prop:curvature_pullback_compatibility}
Suppose $f : X \to Y$ is a map of $S$-schemes where $Y \to S$ is smooth. Suppose there is a commutative diagram
\begin{center}
\begin{tikzcd}
\F_{X} \arrow[d, "\kappa"] \arrow[r, "\alpha"] & \T_{X/S} \arrow[d, "\d{f}"]
\\
\F_{Y} \arrow[r, "f^* \beta"] & f^* \T_{Y/S}
\end{tikzcd}
\end{center}
where $\F_X$ and $\F_Y$ are coherent sheaves on $X$ and $Y$ respectively. As before consider the curvature maps
\[ \ell_{\F_{X}} : \bigwedge\nolimits^{\! 2} \F_X \to \T_{X/S} / \F_X  \quad \text{ via } \quad \ell_{\F_X} : X \wedge Y \mapsto [X,Y] \mod \F \] 
and likewise for $\ell_{\F_{Y}}$. These maps are compatible in the sense that the following diagram relating the curvatures
\begin{center}
\begin{tikzcd}
\bigwedge\nolimits^{\! 2} \F_X \arrow[d, "\wedge^2 \kappa"'] \arrow[r, "\ell_{\F_X}"] & \T_{X/S} / \F_{X} \arrow[d, "\d{f}"]
\\
f^* \bigwedge\nolimits^{\! 2} \F_{Y} \arrow[r, "f^* \ell_{\F_Y}"] & f^* \T_{Y/S} / \F_{Y}
\end{tikzcd}
\end{center}
commutes.
\end{prop}

\begin{proof}
We turn this into the corresponding algebra statement: let $\phi : A \to B$ be a map of $R$-algebras, $N$ an $A$-module and $M$ a $B$-module equipped with a commutative diagram of $B$-modules,
\begin{center}
\begin{tikzcd}[column sep = large]
\Der[R]{B}{B} \arrow[r, "\partial \mapsto \partial \circ \phi"] & \Der[R]{A}{B} & B \ot_A \Der[R]{A}{A} \arrow[l, equals]
\\
M \arrow[u, "\alpha"] \arrow[rr, "\kappa"] & & B \ot_A N \arrow[u, "\id_B \ot \beta"']
\end{tikzcd}
\end{center}
then we need to show that $[\alpha(m), \alpha(m')] = [\beta(\kappa(m)), \beta(\kappa(m'))]$ modulo $M$ where the second Lie bracket is defined by taking the image of $\kappa(m)$ in $B \ot_A \Der[R]{A}{A}$ and extending the bracket $B$-linearly. This is well-defined up to elements of $N \ot_A B$. Explicitly, write,
\[ \beta(\kappa(m)) = \sum f_i \ot \partial_i \]
for some $\partial_i \in \Der[R]{A}{A}$ in the image of $\beta : N \to \Der[R]{A}{A}$ and $f_i \in B$ and likewise
\[ \beta(\kappa(m')) = \sum f'_i \ot \partial_i. \]
Then,
\[ [\beta(\kappa(m)), \beta(\kappa(m'))] = \sum f_i f_j' [\partial_i, \partial_j] \in \Der[R]{A}{B}/(N \ot_A B). \]
However, $\beta(\kappa(m))$ has image $\alpha(m) \circ \phi$ in $\Der[R]{A}{B}$. Therefore, for $a \in A$,
\begin{align*}
[\alpha(m), \alpha(m')](a) &= \alpha(m) \alpha(m') \phi(a) - \alpha(m') \alpha(m) \phi(a) 
\\
& = \alpha(m) \left[ \sum_j f_j' \partial_j a \right] - \alpha(m') \left[ \sum_i f_i \partial_i a \right]
\\
& = \sum_j  \left[ f_j' \alpha(m) (\partial_j a) + \alpha(m)(f_j') \cdot (\partial_j a) \right] - \sum_i \left[ f_i \alpha(m') (\partial_i a) + \alpha(m')(f_i) \cdot (\partial_i a) \right]
\\
& = \sum_{i,j} f_i f_j' [\partial_i, \partial_j](a) + \left[ \sum_j \alpha(m)(f_j') \cdot \partial_j - \sum_i \alpha(m')(f_i) \cdot \partial_j \right]
\end{align*}
but $\partial_i$ is in the image of $N \to \Der[R]{A}{A}$ so the last term is zero in $\Der[R]{A}{B} / B \ot_A N$ proving the claim. 
\end{proof}

Likewise, when $\F$ is not $p$-closed, its $p$-curvature also provides an obstruction to the existence of algebraic leaves. For example, the following elementary observation is quite powerful. 

\begin{prop}
Let $X$ be a smooth variety and $f : Z \to X$ a morphism invariant under a foliation $\F \subset \T_X$ in the sense that the tangent map factors as
\[ \d{f} : \T_Z \to f^* \F \to (f^* \Omega_Y)^{\vee} \]
then the composition of $\T_Z \to f^* \F$ with the pullback of the $p$-curvature
\[ \T_Z \to \F \xrightarrow{f^* \psi_{\F,p}} f^* (\T_X / \F)  \]
is identically zero.   
\end{prop}

\begin{proof}
This follows immediately from the following Proposition~\ref{prop:pullback_p_curvature} expressing the compatibility of $p$-curvature with pullbacks applied to the map $f : (Z, \T_Z) \to (X, \F)$ of foliated varieties noting that the $p$-curvature of the trivial foliation $\T_C \subset \T_C$ is zero. 
\end{proof}

Our primary application of this result is in the case of (families of) curves. We have assumed previously that $f : Z \to X$ is invariant in the strong sense that $\d{f}$ factors globally through the foliation. When we defined a leaf, it was important to relax this condition along the singular locus of the foliation. Similarly, here we only require that the tangent space lies in $\F$ generically. This is only a generalization in the presence of singularities. Indeed, since $\T_Z$ is torsion-free, along the smooth locus of $\F \subset \T_X$ where $\F$ is a sub-bundle, the tangent map $\d{f}$ factors globally if and only if it factors at the generic point. 

\begin{cor}
Let $X$ be a smooth variety and $f : C \to X$ a generically immersed map from a smooth curve generically invariant under a foliation $\F$. Then
\[ f(C) \subset \Sing{\F} \cup \ol{K^1(\psi_{\F,p})} \]
where $K^1$ is the locus where a map of $\cO_X$-modules has nontrivial kernel.
\end{cor}

\begin{proof}
Assume $f(C) \not\subset \Sing(\F)$ then by definition there is a dense open subset $U \subset C$ such that $f(U)$ is contained in the smooth locus of $X$ and $\F$ and there is a factorization
\begin{center}
\begin{tikzcd}
\T_U \arrow[rd] \arrow[r, "\d{f}"] & f|_{U}^* \T_X \arrow[d]
\\
& f|_{U}^* \F \arrow[u, hook]
\end{tikzcd}
\end{center}
where $\d{f}$ is nonvanishing on $U$. By the previous proposition, the composition,
\[ \T_U \to f|_{U}^* \F \xrightarrow{f|_U^* \psi_{\F,p}} f|_{U}^* (\T_X / \F) \]
is identically zero. Since $\T_U  \to f|_{U}^* \F$ is nonvanishing, $f(U) \subset K^1(\psi_{\F,p})$ and therefore $f(C) \subset \ol{K^1(\psi_{\F,p})}$. 
\end{proof}

For example, if $\F$ is a rank $1$ foliation then,
\[ \im{f} \subset \Sing{\F} \cup V(\psi_{\F,p}) \]
If $X$ is a surface and $\F$ is not $p$-closed then $\Sing{\F}$ is a finite collection of points so any invariant curve is contained in $\Delta_{\F,p} := V(\psi_{\F,p})$. In particular, on a surface, there are only finitely many curves tangent to a foliation with nonvanishing $p$-curvature. Now we prove the compatibility of $p$-curvature with pullback.


\begin{prop} \label{prop:pullback_p_curvature}
Let $f : (X, \F_X) \to (Y, \F_Y)$ be a map of foliated $S$-schemes with $Y \to S$ smooth. Explicitly, this means there is a commutative diagram
\begin{center}
\begin{tikzcd}
\F_{X} \arrow[d, "\kappa"] \arrow[r, "\alpha"] & \T_{X/S} \arrow[d, "\d{f}"]
\\
\F_{Y} \arrow[r, "f^* \beta"] & f^* \T_{Y/S}
\end{tikzcd}
\end{center}
and the subsheaves $\F_X$ and $\F_Y$ are assumed to be involutive. Then the diagram
\begin{center}
\begin{tikzcd}
\Frob_X^* \F_X \arrow[r, "\psi_{\F,p}"] \arrow[d] & \T_{X/S} / \F_X \arrow[d]
\\
f^* \Frob_Y^* \F_Y \arrow[r, "f^* \psi_{\F,p}"] & f^* (\T_{Y/S} / \F_Y)
\end{tikzcd}
\end{center}
of $\struct{X}$-linear maps commutes.
\end{prop}

First we need a lemma.


\begin{lemma}
Let $A,B$ be (possibly noncommutative) associative rings of characteristic $p$ and let $M$ be an $(A,B)$-bimodule. Suppose that for some $a \in A$ and $m \in M$,
\[ a \cdot m = \sum_i f_i \cdot m \cdot b_i \]
such that the $f_i \in A$ commute with each other and $b_i \in B$. Then the difference,
\[ a^p \cdot m - \sum_i f_i^p \cdot m \cdot b_i^p \]
lies in the subspace
\[  \left< r \cdot m \cdot s \mid r \in \ZZ[f_1, \dots, f_r, \Lie(a, f_1, \dots, f_r)_{p-1}] \text{ and } s \in \Lie(b_1, \dots, b_r)_{p-1} \right>  \]
where $\Lie(x_1, \dots, x_r)_{p-1}$ denotes the $(p-1)$-th stage of the lower central series of the Lie algebra generated by $x_1, \dots, x_r$. 
\end{lemma}


\begin{proof}
$M$ is an $A \ot B^\op$-module and $(a \ot 1 - \sum_i f_i \ot b_i)$ acts as zero. Furthermore,
\[ (a \ot 1 - \sum_i f_i \ot b_i)^p = a^p \ot 1 + \left(- \sum_i f_i \ot b_i \right)^p + \sum_{j = 1}^{p-1} s_j(a \ot 1, -\sum_i f_i \ot b_i) \]
so the last term is in the span of $\Lie(a, f_1, \dots, f_r)_{p-1} \ot b_i$. By Jacobson again,
\[  \left(- \sum_i f_i \ot b_i \right)^p = - \sum_i f_i^p \ot b_i^p + g \]
where $g$ is in the span of $f_i \ot \Lie(b_1, \dots, b_r)_{p-1}$. Therefore, applying this element to $m \in M$ we obtain,
\[ a^p \cdot m - \sum_i f_i^p \cdot m \cdot b_i^p = -\sum_{j = 1}^{p-1} s_j(a \ot 1, \sum_i f_i \ot b_i) \cdot m  -g \cdot m. \]
\end{proof}



\begin{proof}[Proof of proposition]
This is a local statement so we may reduce to an $R$-algebra map $\phi : A \to B$ with $R \to A$ smooth along with a diagram of $B$-modules
\begin{center}
\begin{tikzcd}[column sep = large]
\Der[R]{B}{B} \arrow[r, "\partial \mapsto \partial \circ \phi"] & \Der[R]{A}{B}
\\
M \arrow[u] \arrow[r, "\kappa"] & B \ot_A N \arrow[u]
\end{tikzcd}
\end{center}
Suppose that $\partial \in M$ we need to show that 
\[ \psi_p(\Frob^* \partial) \circ \phi - \psi_p(\Frob^* \kappa(\partial)) \in \im{( B \ot_A N \to \Der[R]{A}{B})}. \] 
We can expand the $\kappa(\partial) \in B \ot_A N$ 
\[ \kappa(\partial) = \sum_i b_i \ot \partial_i \]
in terms of elements $\partial_i \in N$ and scalars $b_i \in B$. Recall that $\psi_p$ is $\Frob$-semilinear so $f^* \psi_p$ is $B$-semilinear meaning,
\[ \psi_p(\kappa(\partial)) = \psi_p \left(\sum_i b_i \ot \partial_i \right) = \sum_i b_i^p \ot \partial_i^p. \]
Let $\D_A$ be the algebra consisting of differential operators on $A$ generated by derivations. We apply the lemma to $B \ot_A \D_A$ viewed as a $(\D_B, \D_A)$-bimodule. The left $\D_B$ action is defined via
\[ g \cdot (b \ot D) = gb \ot D \quad \text{and} \quad \partial \cdot (b \ot D) = (\partial b) \ot D + b \cdot (\partial \circ \phi) D \]
where, using smoothness of $A$ over $R$, $\partial \circ \phi \in \Der[R]{A}{B} = B \ot_A \Der[R]{A}{A}$ so the element $(\partial \circ \phi) D$ is a well-defined element of $B \ot_A \D_A$. This $(\D_B, \D_A)$-bimodule is sometimes called the ``transfer module'' in $D$-module theory.
Commutativity of the diagram involving $\kappa$ shows that, viewing $\partial$ and $\partial_i$ as differential operators in $\D_B$ and $\D_A$ respectively, we have the relation,
\[ \partial \cdot 1 = \sum_i b_i \cdot 1 \cdot \partial_i \]
in $B \ot_A \D_A$. Therefore, applying the lemma we see that
\[ \psi_p(\partial) \circ \phi - \psi_p(\kappa(\partial)) = \partial^p \circ \phi - \sum_i f_i^p \ot \partial_i^p \in B \ot_A \D_A \]
is in the submodule spanned by elements of the form $r \cdot 1 \cdot s = rs$ where $r$ is in the subalgebra generated by iterated application of $\partial$ to the scalars $b_i \in A$ and products of the $b_i$ and $s$ is generated by commutators of $\partial_i$ which lie in $N$. Hence the difference lies in the image of $B \ot_A N$ completing the proof.
\end{proof}


\section{Ekedahl's theory of positive characteristic foliations} \label{section:Ekedahl}

We now review the theory of $p$-closed foliations, what Ekedahl calls $1$-foliations, as developed by Ekedahl \cite{ekedahl:foliations}. In particular, there is an equivalence between $p$-closed foliations and certain purely inseparable morphisms. This correspondence was first understood by Jacobson in the context of developing a Galois theory for purely inseparable extensions \cite{jacobson:galois_theory_purely_inseparable} and realized in its modern form by Ekedahl in \cite{ekedahl:foliations}. From now on in this section, let $k$ be a perfect field of characteristic $p$ and all varieties will be varieties over $k$. As before, we let $\Frob_X : X \to X$ denote the ($k$-semilinear) absolute Frobenius and $F_{X/k} : X \to X^{(p)}$ denote the ($k$-linear) relative Frobenius.


\begin{defn}
Let $f : X \rat Y$ be a dominant rational map of varieties. We say that $f$ is \textit{purely inseparable} if the extension of function fields $k(Y) \embed k(X)$ is purely inseparable. If, moreover, $f$ is generically finite, then 
$k(X)^{p^h} \subset k(Y)$ for some $h$. The smallest such $h$ is called the \textit{height} of $f$.
Note that if $f$ has height $h$ then by definition there is a rational map $g : Y \rat X^{(p^h)}$ so that $g \circ f : X \rat Y \rat X^{(p^h)}$ is $h$-iterated relative Frobenius $F_{X/k}^h : X \to X^{(p^h)}$ giving a factorization $g \circ f = F_{X/k}^h$.
\end{defn}

\begin{theorem}[Ekedahl] \label{thm:ekedahl}
Let $X$ be a smooth variety. There is an equivalence between the following two posets
\begin{enumerate}
    \item $p$-closed foliations $\F \subset \T_X$
    \item finite height-$1$ purely inseparable morphisms $f : X \to Y$ with $Y$ normal
\end{enumerate}
constructed as follows: given a $p$-closed foliation $\F \subset \T_X$ we define $f : X \to X / \F$ where $f$ is the identity of the underlying topological spaces and $\cO_{X/\F} = \Ann_{\F}(\cO_X)$ where explicitly
\[ \struct{X/\F}(U) := \{ s \in \struct{X}(U) \mid \forall \partial \in \F(U) : \partial s = 0 \}. \]
Therefore, $f$ is the map $X \to \mathbf{Spec}_X(\Ann_{\F}(\cO_X))$.
In reverse, given $f : X \to Y$ the $p$-closed foliation is the subsheaf $\T_{X/Y} \subset \T_X$ of relative tangent fields or equivalently $\Ann_{\T_X}(f^{-1} \cO_Y)$. Furthermore, the correspondence has the following properties:
\begin{enumerate}
    \item $Y$ is smooth iff $\F$ is \textit{nonsingular} ($\F \subset \T_X$ is a subbundle) iff $f$ is flat
    \item over the open set $U := X \sm \Sing(\F)$ where $\F \subset \T_X$ is a sub-bundle, there is an exact sequence of $\cO_X$-modules
    \[ 0 \to \F \to \T_X \to f^* \T_Y \to \Frob_X^* \F \to 0 \]
    which implies that, over $U$,
    \[ K_X = f^* K_Y + (p-1) c_1(\F) \]
    \item over $U$, the map $f^* \T_Y \to \Frob_X^* \F$ is the pullback of a map $\T_Y \to g^* \sigma^* \F$ where $g : Y \to X^{(p)}$ is such that $g \circ f = F_X$ and $\sigma : X^{(p)} \to X$ is the canonical $k$-semilinear isomorphism. Furthermore, the kernel of $\T_Y \to g^* \sigma^* \F$ is a $p$-closed foliation defining $g$. 
\end{enumerate}
\end{theorem}

In general, $Y$ can have quite nasty singularities so $K_Y$ may only be represented by a Weil divisor. In this case, $f^* K_Y$ may not be well-defined except over $U$. We will use Ekedahl's constructions to understand the leaves of a foliation $\F$ in characteristic $p$. When $\F$ is $p$-closed, we get the following useful characterization of the leaves of $\F$.

\begin{lemma} \label{lem:tangent_leaf_pullback}
Let $\F$ be a $p$-closed foliation on a smooth variety $X$. Let $Z \subset X$ be a subvariety of dimension $\rank{\F}$ not contained in the singular locus of $\F$. Then $Z$ arises as the (generically reduced) pullback of a subvariety $Z' \subset X / \F$ if and only if $Z$ is a leaf of $\F$.    
\end{lemma}

\begin{proof}
There is always a commutative diagram of relative Frobenii,
\begin{center}
    \begin{tikzcd}
    Z \arrow[r, "F_Z"] \arrow[d] & Z^{(p)} \arrow[d]
    \\
    X \arrow[r, "F_X"] & X^{(p)}
    \end{tikzcd}
\end{center}
This diagram is not Cartesian: the fiber product is rather the extension of $Z$ by nilpotents along its conormal sheaf. If $Z$ is a leaf, the map $\T_Z \to \T_X|_Z$ is identified with (the reflexive hull) of $\F_X|_Z \to \T_X$ therefore there is a compatible diagram of quotients by foliations
\begin{center}
    \begin{tikzcd}
     Z \arrow[d] \arrow[r] & Z^{(p)} \arrow[d] \arrow[r, equals] & Z^{(p)} \arrow[d]
     \\
     X \arrow[r] & X / \F \arrow[r] & X^{(p)}
    \end{tikzcd}
\end{center}
so the unique map $Z \to Z^{(p)} \times_{X / \F} X$ must be an isomorphism at the generic point by degree considerations. 
\par 
Since $Z$ is not contained in the singular locus of $\F$, we may shrink to assume that $Y = X / \F$ is smooth and $f : X \to X / \F = Y$ is flat. Now, let $Z' \subset X / \F$ be a subvariety and consider the fiber product $F := Z' \times_{X / \F} X$. By base change, the relative cotangent bundle $\Omega_{F/Z'}$ is the pullback of $\Omega_{X/Y} \cong \F^{\vee}$ where $f : X \to X / \F$ is the quotient map. Hence, if $F$ is generically reduced, i.e., the case when $Z$ is the preimage of $Z'$, then $Z \to Z'$ must have zero differential. Therefore, from the diagram
\begin{center}
    \begin{tikzcd}
    & & \T_Z \arrow[r, "0"] \arrow[d] \arrow[ld, dashed] & \T_{Z'}|_Z \arrow[d]
    \\
    0 \arrow[r] & (\F|_Z)^{\vee \vee} \arrow[r] & \T_X |_Z \arrow[r] & \T_{X / \F}|_Z
    \end{tikzcd}
\end{center}
we see that $\T_Z \to \T_X|_Z$ factors through $(\F|_Z)^{\vee \vee} \to \T_X|_Z$. 
\end{proof}




\ifSubfilesClassLoaded{%
  \bibliographystyle{alpha}
  \bibliography{refs}
}{}

\end{document}
